\begin{obereflection}
\textbf{Latihan 8.1}
1. Mengapa tahap \textit{MixColumns} dihilangkan pada putaran terakhir AES?
2. Jelaskan bagaimana kombinasi \textit{SubBytes}, \textit{ShiftRows}, dan \textit{MixColumns} secara kolektif mengimplementasikan prinsip \textit{confusion} dan \textit{diffusion} Shannon.
3. Berikan analisis singkat: mengapa AES-256 lebih tahan terhadap serangan \textit{quantum computing} (misal: algoritma Grover) dibandingkan AES-128?

\textbf{Latihan 8.2}
4. Jelaskan mengapa pemetaan 16 byte masukan ke matriks \textit{state} harus dilakukan
   kolom demi kolom, bukan baris demi baris. Apa akibat praktisnya bagi implementasi
   yang salah memilih konvensi?
5. Buktikan bahwa AddRoundKey bersifat involutif, sedangkan SubBytes dan MixColumns
   tidak. Untuk yang terakhir, tunjukkan dua byte berbeda yang dipertukarkan oleh
   SubBytes.
6. Hitunglah dengan tangan $\code{02} \bullet \code{57}$ dan
   $\code{04} \bullet \code{57}$ di $\text{GF}(2^8)$, lalu tentukan
   $\code{08} \bullet \code{57}$. Petunjuk: setiap langkah hanya menggeser satu bit
   ke kiri dan mereduksi dengan $\code{1B}$ bila perlu.
7. Diketahui bahwa invers perkalian dari $\code{1F}$ di $\text{GF}(2^8)$ adalah
   $\code{B2}$. Periksa dengan Persamaan~\eqref{eq:aes-sbox-rumus} bahwa
   $\text{S-box}(\code{1F}) = \code{C0}$, lalu periksa hasil itu pada
   Tabel~\ref{tab:sbox-aes}.
8. Bangkitkan $\text{RC}[11]$ seandainya AES memerlukan sebelas kunci putaran
   untuk varian berputaran 11. Petunjuk: lanjutkan dari
   $\text{RC}[10] = \code{36}$.

\textbf{Latihan 8.3}
9. Tuliskan urutan transformasi yang dijalankan oleh satu putaran dekripsi AES dan
   bandingkan dengan urutan pada putaran enkripsi. Jelaskan mengapa keduanya tidak
   sekadar kebalikan cermin.
10. Pada \textit{equivalent inverse cipher}, kunci putaran diubah menjadi
    $\tilde{K}_i = \text{InvMixColumns}(K_i)$. Tunjukkan dengan mempergunakan
    sifat linier InvMixColumns bahwa perubahan itu tidak mengubah hasil dekripsi.
11. Tunjukkan bahwa InvShiftRows dan InvSubBytes saling komutatif, tetapi
    InvMixColumns tidak komutatif dengan keduanya. Berikan satu contoh penyangkal
    untuk yang terakhir.
12. Untuk AES-256, jelaskan mengapa aturan $\text{SubWord}$ tambahan diperlukan
    hanya ketika $i \bmod 8 = 4$, dan apa yang terjadi pada jadwal kunci bila aturan
    itu dihilangkan.
13. Berapa banyak word yang harus dibangkitkan untuk AES-192, dan berapa nilai
    $\text{Rcon}$ terbesar yang dipakai? Jelaskan asal angka-angka itu.
14. AES-128 memerlukan 10 putaran, sedangkan DES 16 putaran atas blok yang justru
    lebih kecil. Jelaskan mengapa AES tetap dapat mencapai difusi penuh lebih cepat.
\end{obereflection}
